\documentclass[11pt]{article}
\usepackage[margin=1in]{geometry}
\usepackage{booktabs}
\usepackage{longtable}
\usepackage{array}
\usepackage{hyperref}
\usepackage{xcolor}
\usepackage{enumitem}
\usepackage{parskip}
\usepackage{amssymb}
\hypersetup{colorlinks=true, linkcolor=blue, urlcolor=blue, citecolor=blue}

\title{Replication Report: Singh / Urbaniak et al.\ 2018 \\
\large Multi-drug resistant \textit{Enterobacter bugandensis} from the ISS \\
\normalsize BVBRC-30 $\cdot$ Verdict: \textbf{PARTIAL (strong core replication)}}
\author{Independent Replication (BVBRC-30, CherryRd, local free compute)}
\date{2026-07-01}

\begin{document}
\maketitle

\section*{Paper}
Singh NK, Bezdan D, Checinska Sielaff A, Wheeler K, Mason CE, Venkateswaran K.
``Multi-drug resistant \textit{Enterobacter bugandensis} species isolated from the International Space Station
and comparative genomic analyses with human pathogenic strains.''
\textit{BMC Microbiology} (2018) 18:175. DOI: 10.1186/s12866-018-1325-2. PMID: 30466389. PMCID: PMC6251167.

\section{Executive Summary}
Three independent LLM judges scored the replication as 2$\times$PARTIAL, 1$\times$REPLICATED
(means: Coverage 7.7, Agreement 8.0, Fidelity 7.0, Reproducibility 7.0).
Using the real deposited genome assemblies (BioProject PRJNA319366; exact WGS accessions from the paper's
Table~1), the central computational claims were independently reproduced:
\begin{itemize}[leftmargin=1.3em]
  \item \textbf{Species identity (ANI):} the five ISS isolates are \textit{E.~bugandensis} --- ANI to the three
        clinical comparators reproduced within \textbf{0.03--0.30\%} of the paper's Table 1;
        other \textit{Enterobacter} species fall below the ${\sim}91\%$ species boundary. \checkmark
  \item \textbf{Near-clonality:} all five ISS strains are essentially identical --- fastANI $\geq$99.988\%
        and an identical MLST sequence type \textbf{ST2504}. \checkmark
  \item \textbf{AMR/MDR repertoire:} identical core resistance set --- \textbf{blaACT} (AmpC class-C
        $\beta$-lactamase), \textbf{fosA}, \textbf{oqxA/oqxB} (RND multidrug efflux), \textbf{fieF} (metal efflux) ---
        matching the mechanisms in the paper's Table 2 and explaining the reported
        cephalosporin/cefoxitin phenotype. \checkmark
\end{itemize}
Two items were \textbf{not} computationally reproduced: exact SNP-distance counts (method-dependent)
and wet-lab antibiotic susceptibility phenotypes (require cultures/Vitek).

\section{Paper Claims Tested}
\begin{longtable}{@{}c p{9.5cm} l l@{}}
\toprule
ID & Claim & Type & Status \\ \midrule
C1 & 5 ISS isolates are \textit{E.~bugandensis}: ANI ${>}95\%$ to clinical strains, ${<}91\%$ to other species; dDDH ${\sim}89\%$ to EB-247/153\_ECLO & ANI/dDDH & Replicated \\
C2 & 5 ISS strains near-clonal (max 15 SNPs among them) & clonality & Replicated (qualitatively) \\
C3 & Broad AMR/MDR gene repertoire ($\beta$-lactamase, RND efflux, MAR, metal) & AMR detection & Replicated \\
C4 & Phenotypic resistance to cefazolin, cefoxitin, oxacillin, penicillin, rifampin & wet-lab phenotype & Not reproducible (genotype consistent) \\
C5 & ${\sim}4733$ genes IF3SW-P2; carbohydrate/AA/protein-metabolism dominant; 112 virulence genes & annotation counts & Consistent (RAST not re-run) \\
C6 & 16S ${\sim}99.6\%$ to EB-247T; MLST/gyrB places ISS with EB-247/153\_ECLO & phylogenetic placement & Superseded by MLST+ANI \\
\bottomrule
\end{longtable}

\section{Methods \& Data Provenance}
Genomes downloaded via NCBI Datasets 18.25.1 on 2026-07-01: 5 ISS isolates (IF3SW-P2, IF2SW-P2,
IF2SW-B1, IF2SW-P3, IF2SW-B5), 3 clinical \textit{E.~bugandensis} comparators (EB-247T, 153\_ECLO, MBRL1077),
and 5 outgroup \textit{Enterobacter} type/reference strains (\textit{E.~cloacae} ATCC13047, \textit{E.~asburiae}
ATCC35953, \textit{E.~ludwigii} EN-119, \textit{E.~aerogenes} KCTC2190, \textit{E.~kobei}).
ISS assemblies ${\sim}4.93$~Mb, ${\sim}55.9\%$ GC, 2 contigs (hybrid Nanopore+Illumina).

Tools (all free/local): fastANI (all-vs-all ANI; paper used the Goris 2007 BLAST-ANI algorithm),
AMRFinderPlus 4.2.7 (AMR genes; paper used RAST subsystems), mlst 2.33.1 (\textit{E.~cloacae}-complex
scheme, same DB the paper used), minimap2 (asm5) + paftools (SNP calling), biopython 1.87 for stats.

\section{Results}

\subsection{C1 --- Species identity by ANI (headline claim)}
\begin{longtable}{@{}l r r r@{}}
\toprule
Comparator & Replicated ANI\% & Paper ANI\% & $\Delta$ \\ \midrule
\textit{E.~bugandensis} EB-247T & 98.63 & 98.66 & $-0.03$ \\
\textit{E.~bugandensis} 153\_ECLO & 98.64 & 98.73 & $-0.09$ \\
\textit{E.~bugandensis} MBRL1077 & 95.56 & 95.26 & $+0.30$ \\
\textit{E.~kobei}                 & 91.11 & 90.54 & $+0.57$ \\
\textit{E.~ludwigii} EN-119       & 88.40 & 87.57 & $+0.83$ \\
\textit{E.~cloacae} ATCC13047     & 88.84 & 87.91 & $+0.93$ \\
\textit{E.~asburiae} ATCC35953    & 91.89 & 85.59 & $+6.30$ \\
\textit{E.~aerogenes} KCTC2190    & 81.91 & 78.74 & $+3.17$ \\
\bottomrule
\end{longtable}
The three \textit{E.~bugandensis} clinical strains reproduce essentially exactly
($\Delta \leq 0.30\%$). Larger deviations for \textit{E.~asburiae}/\textit{E.~aerogenes} reflect
different downstream assemblies and fastANI vs Goris BLAST-ANI diverging at lower identity.
The species-defining conclusion --- all 5 ISS strains are \textit{E.~bugandensis} ---
is fully reproduced.

\subsection{C2 --- Near-clonality}
\begin{itemize}[leftmargin=1.3em]
  \item fastANI ISS-vs-ISS: minimum 99.988\%, most pairs 99.999\%.
  \item MLST: all five ISS strains = identical \textbf{ST2504}; clinical strains distinct
        (EB-247=ST495, 153\_ECLO=ST659). Independent confirmation of clonality not present
        in the original paper.
  \item SNP counts (assembly-vs-assembly, minimap2+paftools) vs IF3SW-P2: \textbf{81--183 SNPs} ---
        higher than the paper's \textbf{9--15} (bwa-mem read mapping + GATK HaplotypeCaller with
        false-positive filtering). Even at the stricter count, identity is ${>}99.996\%$.
        Qualitative clonality claim replicated; exact SNP numbers method-dependent.
\end{itemize}

\subsection{C3 --- AMR / MDR repertoire}
All 5 ISS strains share an \textbf{identical} AMRFinderPlus core AMR set:
\begin{itemize}[leftmargin=1.3em]
  \item \texttt{blaACT} --- AmpC class-C $\beta$-lactamase (paper Table 2: ``Beta-lactamase class C and other PBPs'')
  \item \texttt{fosA}   --- fosfomycin resistance
  \item \texttt{oqxA} / \texttt{oqxB} --- RND multidrug efflux pump (tripartite MDR / RND efflux system)
  \item \texttt{fieF}   --- ferrous-iron / metal (Co/Zn/Cd) efflux
\end{itemize}
The AmpC $\beta$-lactamase (\texttt{blaACT}) is mechanistically consistent with the paper's reported
phenotypic resistance to \textbf{cefazolin} and \textbf{cefoxitin} (C4). Clinical strains carry a broader
set (\texttt{silA}, \texttt{qnrE}, \texttt{blaIMI-1} carbapenemase in MBRL1077, extra \texttt{fosA7}),
consistent with the paper's theme of expanded AMR in clinical isolates.

\subsection{C5 --- Genome statistics}
ISS genomes ${\sim}4.93$ Mb, ${\sim}55.9\%$ GC, 2 contigs --- internally consistent with the paper's
${\sim}4733$-gene report for IF3SW-P2 (${\sim}1$ gene/kb) and the \textit{Enterobacter} GC range.
RAST subsystem gene-category counts (635/496/291\ldots) were not regenerated
(would require a RAST run); marked consistent, not re-derived.

\section{Discrepancies \& Limitations}
\begin{enumerate}[leftmargin=1.5em]
  \item \textbf{SNP counts (9--15 paper vs 81--183 here)} --- assembly-vs-assembly vs filtered read-mapping.
        Different methods, same clonality conclusion.
  \item \textbf{Outgroup ANI deviations} for \textit{E.~asburiae}/\textit{E.~aerogenes} ($\Delta$ 3--6\%)
        --- different reference assemblies + fastANI vs Goris ANI at lower identity.
        Does not affect species-level conclusions.
  \item \textbf{Wet-lab phenotypes (C4)} --- disk diffusion/Vitek not reproducible in silico;
        AMR genotype is consistent with them.
  \item \textbf{RAST subsystem counts (C5)} --- not regenerated; AMRFinderPlus (different paradigm)
        used for AMR, with concordant core genes.
  \item \textbf{dDDH} --- not recomputed (paper used GGDC web service); ANI serves as the equivalent
        species-boundary metric.
\end{enumerate}

\section{Reproducibility Assessment}
High. All genomes are public with verified accessions; all tools are free and versioned; analysis
runs on a laptop in minutes (5~Mb bacterial genomes). The paper's precise SNP pipeline
(raw reads + GATK filters) and RAST annotation counts are the only pieces not rerunnable from
these artifacts.

\section{Genuine Critique}
This section states, without soft-pedaling, what this replication does \emph{and does not} settle.

\subsection*{What genuinely replicates}
\begin{itemize}[leftmargin=1.3em]
  \item \textbf{Species assignment is rock-solid.} Three independent \textit{E.~bugandensis} comparators
        all land at $98.6$--$95.6\%$ ANI to the ISS quintet, and every non-bugandensis \textit{Enterobacter}
        falls below the $\sim$95\% species boundary. This is the paper's headline claim and it
        reproduces to within 0.03--0.30\% on Table~1's three key rows.
  \item \textbf{Clonality is over-determined.} fastANI $\geq$99.988\% among the 5 ISS isolates
        \emph{and} identical MLST ST2504 is an even stronger clonality signal than the paper
        presented; the paper never reported MLST for these isolates, so the ST identity is a
        genuine independent finding, not a retread.
  \item \textbf{Core AMR mechanisms are the same.} \texttt{blaACT} (AmpC) + \texttt{fosA} +
        \texttt{oqxA/oqxB} + \texttt{fieF} is a coherent chromosomal-AMR profile and it is
        \emph{identical across all 5 ISS strains}, which is itself confirmatory of clonality.
\end{itemize}

\subsection*{What does not replicate, and why the honest verdict is PARTIAL rather than REPLICATED}
\begin{itemize}[leftmargin=1.3em]
  \item \textbf{The SNP-distance number is off by roughly an order of magnitude
        (81--183 vs 9--15).} The delta is fully explained by methodology (assembly-vs-assembly
        with minimap2+paftools inherits assembly ambiguities that stringent read-mapping +
        GATK false-positive filtering removes), and the \emph{clonality conclusion} is unaffected
        (identity still ${>}99.996\%$). But to actually reproduce the paper's SNP number we
        would need the raw Illumina reads and the paper's bwa-mem + GATK HaplotypeCaller
        filter pipeline. Until that is rerun, this is a methodological match, not a numerical
        match, and the report treats it as such (marked $\checkmark/\!\!\bigtriangleup$).
  \item \textbf{Wet-lab phenotype (C4) is by construction out of scope} for a computational
        replication. The AMR genotype (AmpC + fosA + oqxAB) is \emph{mechanistically consistent}
        with cefazolin/cefoxitin/oxacillin/penicillin resistance, but ``consistent with'' is not
        ``reproduces''. Disk diffusion / Vitek data cannot be recovered from assemblies.
  \item \textbf{RAST subsystem counts (C5) were not regenerated.} The paper reports ${\sim}4733$ genes
        and specific subsystem category counts (635/496/291\ldots); this replication used
        AMRFinderPlus for AMR and only sanity-checked genome length/GC/contigs. The gene-total
        and per-subsystem numbers are ``consistent with'' the \textit{Enterobacter} genome
        envelope but not independently re-derived.
  \item \textbf{Outgroup ANI has real numerical drift.} \textit{E.~asburiae} deviates by $+6.30\%$
        and \textit{E.~aerogenes} by $+3.17\%$. These do not affect the species call for the ISS
        isolates, but they \emph{are} real numerical disagreements with Table~1 and reflect the
        (predictable) divergence between fastANI and Goris BLAST-ANI at lower identities plus
        the use of different reference assemblies for the outgroups. If the exercise were
        ``recreate Table~1 to 2 decimals across all rows'', it would fail on those rows.
  \item \textbf{dDDH was not recomputed} at all. The paper's ${\sim}89\%$ dDDH to EB-247/153\_ECLO
        is simply asserted-by-analogy from the ANI values via ``ANI serves as the equivalent
        species-boundary metric''. A strict replication would rerun GGDC.
\end{itemize}

\subsection*{Net honest assessment}
The paper's \emph{scientific} conclusions --- ISS isolates are clonal \textit{E.~bugandensis} carrying
a conserved multidrug-resistance repertoire dominated by chromosomal AmpC, fosfomycin, and RND
efflux --- are independently confirmed on real deposited data with free tooling. The paper's
\emph{numerical table entries} reproduce for the species-defining rows and diverge measurably for
the outgroups and SNP counts. That is exactly what ``PARTIAL (strong core replication)'' means, and
inflating it to ``REPLICATED'' would misrepresent the SNP, RAST, dDDH, and wet-lab gaps that were
not closed.

\section{Artifacts}
\texttt{data/claims.json} (extracted claims + accession map);
\texttt{work/genomes/*.fna} (13 real genome assemblies);
\texttt{work/ani\_matrix.tsv}, \texttt{work/ani\_summary.json} (fastANI results);
\texttt{work/amr/*.tsv}, \texttt{work/amr\_summary.json} (AMRFinderPlus per-strain);
\texttt{work/genome\_stats.json} (length/GC/contigs);
\texttt{work/snp2/*.var} (minimap2+paftools SNP calls);
\texttt{work/analysis\_summary.md} (consolidated evidence);
\texttt{work/judge\_scores.json} (3-judge LLM rubric scores);
\texttt{paper/urbaniak2018.pdf}, \texttt{paper/paper\_extracted.txt} (source).

\end{document}
